\documentclass[11pt]{amsart}

\usepackage{amsmath,amssymb}
\usepackage{hyperref}

\setlength{\emergencystretch}{2em}

\newtheorem{theorem}{Theorem}[section]
\newtheorem{proposition}[theorem]{Proposition}
\newtheorem{corollary}[theorem]{Corollary}
\newtheorem{lemma}[theorem]{Lemma}
\newtheorem{definition}[theorem]{Definition}
\newtheorem{remark}[theorem]{Remark}
\newtheorem{gap}[theorem]{Research gap}

\newcommand{\C}{\mathbb C}
\newcommand{\R}{\mathbb R}
\newcommand{\Hilb}{\mathcal H}
\newcommand{\Span}{\operatorname{span}}
\newcommand{\rank}{\operatorname{rank}}
\newcommand{\Tr}{\operatorname{Tr}}
\newcommand{\Ran}{\operatorname{Ran}}
\newcommand{\Ker}{\operatorname{Ker}}
\newcommand{\Id}{I}
\newcommand{\Proj}{\Pi}
\newcommand{\dd}{\mathrm d}
\newcommand{\FS}{\mathrm{FS}}
\newcommand{\QGT}{\mathrm{QGT}}
\newcommand{\Gr}{\mathrm{Gr}}
\newcommand{\diag}{\operatorname{diag}}

\title[Jet-Gram structures and higher quantum geometry]{Jet-Gram structures,
canonical subspace invariants, and higher quantum geometry}

\author{Draft}

\date{\today}

\begin{document}

\begin{abstract}
The quantum geometric tensor is the first-order Gram matrix of horizontal
derivatives of a pure-state family. Motivated by the local jet cross-Gram
pattern in the Riemann-side construction, we ask whether this first-order
geometric object has a higher-order, two-point matrix analogue for quantum
state families. The answer is sharply mixed. At first order the same-point
horizontal Gram matrix is exactly the usual quantum geometric tensor, and a
transported two-point first-jet block has a precise reparameterization law. From
second order onward, however, ordered jet matrices transform by non-unitary
triangular laws under phase gauge and coordinate change. Their entries are not
canonical. The surviving objects are the horizontal osculating subspaces
\(O_r\), together with the ambient extensions
\(A_r=\Span\{\psi\}\oplus O_r\). Endpoint comparison is then governed by
cross-contractions, projector compressions, polar partial isometries, and
principal angles. We prove that all orthonormal-frame comparison matrices for a
fixed endpoint subspace pair form one biunitary orbit, classified by the same
principal-angle singular values; a canonical entrywise matrix therefore
requires extra frame-selecting structure. We give polynomial and non-polynomial
benchmarks where the resulting higher-order subspace data is richer than
ordinary overlap, record a finite-jet genericity statement for the first
horizontal line invariant, and describe a coherent reflection protocol that
reads the principal angles under explicit access assumptions. Research gaps are
isolated where stronger physics-facing claims would require new input.
\end{abstract}

\maketitle

\tableofcontents

\section{Introduction}

The usual quantum geometric tensor is a first-order object. Given a smooth
family of normalized pure states, one removes the component of a derivative in
the state direction and forms the Gram matrix of the resulting horizontal
derivatives. The real part is the quantum metric; the imaginary part is Berry
curvature. This paper asks what survives when the same Gramian idea is pushed
to higher jets and to two endpoint comparisons.

The motivating algebraic pattern is a local jet cross-Gram block. One has
same-point Gram matrices at two nearby points, a mixed cross block between the
two local jet frames, and a bilateral whitening. In abstract notation this has
the form
\[
        \Omega=M_-^{-1/2}N_{-+}M_+^{-1/2}.
\]
The question is not whether a particular kernel-theoretic object is literally a
standard quantum tensor. The question is narrower and more structural:

\begin{quote}
Which parts of a higher-jet Gram construction are canonical for quantum state
families, and which parts depend on phase, coordinates, frames, or transport?
\end{quote}

The answer is the main message of the paper. At first order one obtains the
standard quantum geometric tensor and a controlled two-point transported block.
At higher order the matrix entries do not survive. The ordered derivative list
mixes by upper-triangular, non-unitary transformations under ordinary changes
of phase and parameter. Whitening does not repair this. The canonical survivor
is the subspace spanned by the derivatives. Comparing endpoints means comparing
subspaces.

Thus the central object is not a higher quantum geometric matrix, but the
endpoint operator package of two osculating subspaces:
\[
        C=\Proj_-|_{S_+},\qquad K_+=C^*C,\qquad
        K_-=CC^*,\qquad C=V|C|.
\]
Its spectra give the squared cosines of principal angles. The polar factor
gives a canonical partial isometry on the nonzero-overlap sector. The
principal-angle multiset classifies finite-dimensional subspace pairs up to
ambient unitary equivalence.

This package has two uses. First, it gives a higher-order geometric fingerprint
of a state family. In explicit qutrit, quartit, Veronese, and spherical-twist
families it distinguishes endpoint pairs with the same ordinary overlap.
Second, it gives a no-go criterion. A proposed higher-jet comparison matrix is
not canonical unless the extra structure selecting its frames is stated.

\subsection{Main contributions}

The paper proves or records the following claims.

\begin{enumerate}
\item The same-point first horizontal Gram matrix is exactly the pure-state
quantum geometric tensor.
\item A one-parameter first-horizontal two-point block is invariant under
orientation-preserving reparameterization after a fixed gauge-covariant
transport choice, and transforms by derivative-channel sign conjugation under
orientation reversal.
\item Raw higher projected jets transform by upper-triangular laws under phase
gauge and reparameterization. Therefore their Gram and whitened matrices are
not canonical.
\item The horizontal osculating spaces \(O_r\) and the ambient spaces \(A_r\)
are canonical subspace-valued invariants.
\item Berry-covariant one- and multi-parameter jet hierarchies generate the
same \(O_r\) filtrations as raw projected jets. Symmetrized covariant jets
generate the same multiparameter filtration.
\item Endpoint comparison is controlled by cross-contractions, projector
compressions, polar partial isometries, and principal angles.
\item Every orthonormal-frame comparison matrix for a fixed subspace pair lies
in a single biunitary orbit; the rectangular singular values are exactly the
principal-angle cosines.
\item Canonical principal frames exist from the subspace pair alone exactly on
simple nonzero principal spectra. Repeated principal values have an unavoidable
internal unitary freedom.
\item Exact unitary orbits identify the ambient and horizontal spaces with
transported Krylov and compressed-Krylov spaces. Kato transport gives a
pathwise realization but not endpoint-independent transport in nonzero
curvature.
\item With coherent reflection access to endpoint projectors, phase estimation
on \((2\Proj_- - \Id)(2\Proj_+ - \Id)\) reads the same principal angles.
\item Benchmark families show richer-than-overlap behavior. A finite-jet
genericity criterion gives a first-order mechanism behind this behavior.
\end{enumerate}

\subsection{What is not claimed}

The paper does not claim the discovery of new fundamental particles. It does
not claim that higher jets define a canonical matrix. It does not claim that the
construction is fundamental for all quantum mechanics. It proves a mathematical
invariant package and marks the additional work needed to turn that package
into a physics-facing diagnostic in concrete systems.

\section{Background and notation}

Let \(\Hilb\) be a complex Hilbert space. Most statements below are cleanest in
finite dimension; the subspace and projector statements extend to closed
finite-rank subspaces in a Hilbert space with the usual domain qualifications.
The inner product is linear in the second variable only if a convention is
chosen; the formulas below are invariant after conjugating consistently. We
write \(|\psi\rangle\) only in informal passages and use \(\psi\) in formulas.

A smooth pure-state family is a smooth map \(x\mapsto \psi(x)\in \Hilb\) with
\(\|\psi(x)\|=1\). The physical state is the ray \([\psi(x)]\), so the local
lift is defined only up to a phase \(e^{i\chi(x)}\). The orthogonal projector
onto the horizontal space at \(\psi\) is
\[
        P_\psi=\Id-|\psi\rangle\langle \psi|.
\]
For a coordinate vector \(\partial_a\) on parameter space, define the horizontal
derivative
\[
        D_a\psi=P_\psi\partial_a\psi .
\]

\begin{definition}[Quantum geometric tensor]
The pure-state quantum geometric tensor of the family \(x\mapsto \psi(x)\) is
\[
        Q_{ab}(x)=\langle D_a\psi(x),D_b\psi(x)\rangle
        =\langle\partial_a\psi,
          (\Id-|\psi\rangle\langle\psi|)\partial_b\psi\rangle .
\]
Its real part is the Fubini-Study metric tensor and twice its imaginary part is
the Berry curvature in the convention used here:
\[
        g_{ab}=\Re Q_{ab},\qquad F_{ab}=2\Im Q_{ab}.
\]
\end{definition}

We also need the elementary geometry of two subspaces. If \(S_-\) and \(S_+\)
are finite-dimensional subspaces with orthogonal projectors \(\Proj_-\) and
\(\Proj_+\), the canonical cross-contraction is
\[
        C=\Proj_-|_{S_+}:S_+\to S_-.
\]
The positive compressions are
\[
        K_+=C^*C=\Proj_+\Proj_-\Proj_+|_{S_+},
        \qquad
        K_-=CC^*=\Proj_-\Proj_+\Proj_-|_{S_-}.
\]
The singular values of \(C\) are the cosines of the principal angles between
\(S_-\) and \(S_+\).

\section{First-order quantum jet-Gram geometry}

\subsection{The same-point bridge}

The first-order same-point part of the construction is exactly standard quantum
geometry.

\begin{theorem}[First-order QGT bridge]
Let \(x\mapsto \psi(x)\) be a smooth normalized pure-state family. The
same-point horizontal derivative Gram matrix
\[
        \bigl(\langle P_\psi\partial_a\psi,
                    P_\psi\partial_b\psi\rangle\bigr)_{a,b}
\]
is the pure-state quantum geometric tensor.
\end{theorem}

\begin{proof}
The statement is the definition after horizontal projection. Since
\(\|\psi\|^2=1\), the derivative \(\partial_a\psi\) decomposes into the state
line and its orthogonal complement. Removing the state-line component gives
\[
        D_a\psi=(\Id-|\psi\rangle\langle\psi|)\partial_a\psi.
\]
Taking the Gram matrix of these horizontal derivatives gives
\[
        \langle D_a\psi,D_b\psi\rangle
        =
        \langle\partial_a\psi,
        (\Id-|\psi\rangle\langle\psi|)\partial_b\psi\rangle,
\]
which is the usual expression for the pure-state quantum geometric tensor.
\end{proof}

This theorem is the exact first-order contact with established physics. It does
not identify the full two-point whitened block with the quantum geometric
tensor; it identifies the same-point horizontal block.

\subsection{A transported first-order two-point block}

Let \(t\mapsto \psi(t)\) be a one-parameter normalized lift. Define
\[
        h_t=P_t\dot\psi(t),\qquad P_t=\Id-|\psi(t)\rangle\langle\psi(t)|.
\]
Choose a two-point transport rule \(U_{-\leftarrow +}\) from the endpoint
fiber at \(t_+\) to the endpoint fiber at \(t_-\). In the ordered pair
\((\psi,h)\), define
\[
M(t_i)=
\begin{pmatrix}
1&0\\
0&\langle h_{t_i},h_{t_i}\rangle
\end{pmatrix}
\]
and
\[
N(t_-,t_+)=
\begin{pmatrix}
\langle\psi_-,U_{-\leftarrow +}\psi_+\rangle&
\langle\psi_-,U_{-\leftarrow +}h_+\rangle\\
\langle h_-,U_{-\leftarrow +}\psi_+\rangle&
\langle h_-,U_{-\leftarrow +}h_+\rangle
\end{pmatrix}.
\]
The whitened first-order block is
\[
        \Omega_t=M(t_-)^{-1/2}N(t_-,t_+)M(t_+)^{-1/2}.
\]

\begin{proposition}[First-order reparameterization law]
Let \(\mu=\phi(t)\) be a smooth reparameterization with nonzero derivative at
the endpoints. For a fixed transported endpoint comparison as above,
\[
        \Omega_\mu=\Omega_t
\]
when \(\phi\) is orientation-preserving at both endpoints. If the orientation is
reversed at both endpoints and the ordered endpoint pair is kept fixed, then
\[
        \Omega_\mu=E\Omega_tE,\qquad E=\diag(1,-1).
\]
\end{proposition}

\begin{proof}
Under reparameterization,
\[
        h_\mu=P_\psi\frac{\partial\psi}{\partial\mu}
        =
        \frac{\dd t}{\dd\mu}h_t.
\]
The derivative channel at each endpoint is therefore multiplied by the real
scalar \(\dd t/\dd\mu\). The corresponding same-point derivative norm in
\(M(t_i)\) is multiplied by its square, while the mixed entries involving one
derivative channel are multiplied by the endpoint derivative factor. Bilateral
whitening cancels the positive magnitudes. Only the signs remain. If both signs
are positive, nothing changes. If both endpoint orientations reverse, the
derivative channel is conjugated by the diagonal sign matrix
\(E=\diag(1,-1)\).
\end{proof}

\begin{remark}[Transport is part of the data]
This proposition is not an endpoint-only theorem. The transport rule is part of
the input. If the transport is gauge-covariant, the construction descends from
a fixed lift to the corresponding ray curve. In the restricted scalar-phase
class \(U_{-\leftarrow +}=e^{i\alpha}\Id\), the whole block changes by a global
phase and its singular values are transport-class invariant. Outside such
restricted classes, no transport-independent matrix is claimed.
\end{remark}

\begin{gap}[GAP-3A: first-order transport classification]
The current draft uses the fixed-transport and scalar-phase statements because
they are enough for the main message. A stronger version would classify all
transport choices under which the first-order singular-value package is
invariant. This is not needed for the paper's core theorem spine.
\end{gap}

\section{Higher jets: matrix failure and subspace survival}

The first-order matrix behavior is exceptional. Higher derivatives mix under
ordinary changes of phase and parameter.

\subsection{Triangular transformation laws}

For a one-parameter curve define raw projected jets
\[
        j_k(t)=P_t\partial_t^k\psi(t),\qquad 1\le k\le r.
\]

\begin{lemma}[Phase triangularity]
Under a phase change \(\psi(t)\mapsto e^{i\chi(t)}\psi(t)\), the ordered list
\((j_1,\ldots,j_r)\) is transformed by multiplication by \(e^{i\chi(t)}\) and
an invertible upper-triangular linear combination of the lower jets. In
particular,
\[
        j_2\mapsto e^{i\chi}(j_2+2i\chi'j_1).
\]
\end{lemma}

\begin{proof}
Differentiate \(e^{i\chi}\psi\) repeatedly. The \(k\)-th derivative is
\[
        \partial_t^k(e^{i\chi}\psi)
        =
        e^{i\chi}\partial_t^k\psi
        +\sum_{\ell<k} b_{k\ell}(\chi',\ldots,\chi^{(k-\ell)})
             e^{i\chi}\partial_t^\ell\psi
        +b_{k0}e^{i\chi}\psi .
\]
The final state-line term is killed by the horizontal projector. The coefficient
of \(j_k\) is \(e^{i\chi}\), so the transformation is triangular with invertible
diagonal. The displayed \(k=2\) formula follows by direct differentiation and
projection.
\end{proof}

\begin{lemma}[Coordinate triangularity]
Let \(\mu=\phi(t)\) be a one-parameter coordinate change with
\(\dd t/\dd\mu\ne0\). Then the ordered list \((j_1,\ldots,j_r)\) transforms by
an invertible upper-triangular combination. In particular,
\[
        j_2^{(\mu)}
        =
        \left(\frac{\dd t}{\dd\mu}\right)^2j_2^{(t)}
        +
        \frac{\dd^2 t}{\dd\mu^2}j_1^{(t)} .
\]
\end{lemma}

\begin{proof}
This is the chain rule. The \(k\)-th \(\mu\)-derivative contains
\((\dd t/\dd\mu)^k\partial_t^k\psi\) plus lower \(t\)-derivatives. Projecting
horizontally kills only the state-line term, not the lower horizontal jets. The
diagonal coefficients are nonzero, so the transformation is invertible and
triangular.
\end{proof}

\begin{theorem}[Higher-jet matrix no-go]
For \(r\ge2\), the raw higher-jet Gram matrices and their bilaterally whitened
two-point matrix analogues are not canonical under phase gauge and
reparameterization. Their entries depend on the ordered jet representatives.
\end{theorem}

\begin{proof}
The transformation laws above are generally non-unitary triangular changes of
ordered jet coordinates. A Gram matrix transforms by congruence under such a
change, and bilateral whitening only normalizes the endpoint same-point Gram
matrices. It does not make the mixed block invariant under arbitrary non-unitary
triangular coordinate changes. Hence the entries of the ordered matrices depend
on gauge and coordinate choices.
\end{proof}

\subsection{Subspace survival}

\begin{theorem}[Osculating subspace survival]
For a one-parameter ray curve, the subspace
\[
        O_r(t)=\Span\{j_1(t),\ldots,j_r(t)\}
\]
is invariant under local phase gauge and under reparameterization with nonzero
derivative. In a fixed ambient Hilbert space the ambient space
\[
        A_r(t)=\Span\{\psi(t),\partial_t\psi(t),\ldots,\partial_t^r\psi(t)\}
              =\Span\{\psi(t)\}\oplus O_r(t)
\]
is also invariant.
\end{theorem}

\begin{proof}
The ordered jets transform by an invertible triangular matrix. Such a matrix
changes the generating list but not its span. Adding the state line gives the
ambient space \(A_r\). The equality
\(A_r=\Span\{\psi\}\oplus O_r\) follows because
\(\partial_t^k\psi=j_k+\langle\psi,\partial_t^k\psi\rangle\psi\).
\end{proof}

This theorem is the pivot of the paper. The failed matrix is replaced by a
canonical subspace.

\subsection{Covariant jets}

The subspace theorem can be generated by covariant derivatives rather than raw
projected coordinate derivatives. In one parameter define
\[
        a(t)=\langle\psi(t),\dot\psi(t)\rangle,\qquad
        \nabla_t\xi=P_t(\dot\xi-a(t)\xi).
\]
Set
\[
        h_{[1]}=\nabla_t\psi,\qquad h_{[k+1]}=\nabla_t h_{[k]}.
\]

\begin{proposition}[One-parameter covariant hierarchy]
The vectors \(h_{[k]}\) are gauge-covariant, and
\[
        \Span\{h_{[1]},\ldots,h_{[r]}\}=O_r(t).
\]
\end{proposition}

\begin{proof}
Under \(\psi\mapsto e^{i\chi}\psi\), a horizontal vector
\(\xi\mapsto e^{i\chi}\xi\) is sent by \(\nabla_t\) to
\(e^{i\chi}\nabla_t\xi\). Thus the hierarchy is gauge-covariant. Expanding
\(\nabla_t h_{[k]}\) shows that \(h_{[k+1]}\) equals \(j_{k+1}\) plus a
linear combination of lower \(j_\ell\)'s. The change of generators is again
triangular with nonzero diagonal, so the spans agree.
\end{proof}

In several parameters, write
\[
        a_a=\langle\psi,\partial_a\psi\rangle,\qquad
        \nabla_a\xi=P(\partial_a\xi-a_a\xi).
\]
Ordered iterated covariant derivatives
\(\nabla_{a_k}\cdots\nabla_{a_1}\psi\) generate the same \(O_r\) spaces as raw
mixed projected partial derivatives of total order at most \(r\). Since
commutators of the \(\nabla_a\)'s contribute only lower-order terms on the
horizontal bundle, symmetrized covariant derivatives generate the same
filtration. Thus the multiparameter hierarchy determines filtered subspaces,
not a canonical ordered frame.

\begin{gap}[GAP-4A: associated-graded symbols]
The multiparameter no-go leaves room for a basis-free associated-graded or
tensor-symbol package. Such an object would not be an endpoint comparison
matrix, but it might be a useful local invariant. A future version should either
define this object and prove its transformation law, or explicitly leave it out
of scope.
\end{gap}

\section{Endpoint operator package}

Once the canonical objects are endpoint subspaces, the comparison problem is
classical and rigid.

\subsection{Canonical contractions}

\begin{theorem}[Endpoint operator package]
Let \(S_-\) and \(S_+\) be finite-dimensional subspaces of \(\Hilb\), with
orthogonal projectors \(\Proj_-\) and \(\Proj_+\). The following data are
canonical:
\[
        C=\Proj_-|_{S_+}:S_+\to S_-,
\]
\[
        K_+=C^*C=\Proj_+\Proj_-\Proj_+|_{S_+},
        \qquad
        K_-=CC^*=\Proj_-\Proj_+\Proj_-|_{S_-},
\]
and the polar decomposition
\[
        C=V|C|
\]
on the nonzero-overlap sector. The nonzero spectra of \(K_+\) and \(K_-\) are
the squared cosines of the principal angles between \(S_-\) and \(S_+\).
\end{theorem}

\begin{proof}
The projectors are determined by the subspaces, so \(C\), \(K_+\), and \(K_-\)
are determined by the subspace pair. The nonzero singular values of \(C\) are
the square roots of the nonzero eigenvalues of \(C^*C\) and \(CC^*\). By the
standard definition of principal angles, these singular values are
\(\cos\theta_i\). The polar decomposition of a bounded finite-rank operator is
unique on \((\Ker C)^\perp\), so the partial isometry \(V\) is canonical on the
nonzero-overlap sector.
\end{proof}

\begin{corollary}[Projector test]
For a state prepared in \(S_-\), the operator
\[
        \Proj_-\Proj_+\Proj_-|_{S_-}
\]
is the effect operator for the sequential yes/yes subspace-membership test. Its
eigenvalues are \(\cos^2\theta_i\).
\end{corollary}

\subsection{Frame matrices are orbit representatives}

The operator package is often written as a matrix after choosing endpoint
frames. The next theorem says exactly what that matrix means.

\begin{theorem}[Biunitary matrix finality]
Let \(U_-:\C^p\to S_-\) and \(U_+:\C^q\to S_+\) be orthonormal frame
isometries. The comparison matrix
\[
        A=U_-^*U_+
\]
is the matrix of \(C=\Proj_-|_{S_+}\) in those frames. Replacing the frames by
\(U_-L\) and \(U_+R\), with \(L\in U(p)\) and \(R\in U(q)\), changes the
matrix by
\[
        A\mapsto L^*AR.
\]
Conversely, every orthonormal-frame comparison matrix for the same pair arises
in this way. Therefore the singular values of \(A\), equivalently the
principal-angle cosines, are the complete frame-independent content of such a
matrix.
\end{theorem}

\begin{proof}
Since \(U_-\) is an isometry onto \(S_-\), \(U_-U_-^*=\Proj_-\) on the ambient
space. For \(z\in\C^q\),
\[
        U_-^*C(U_+z)=U_-^*\Proj_-U_+z=U_-^*U_+z,
\]
so \(A\) is the frame matrix of \(C\). If the frames are replaced by \(U_-L\)
and \(U_+R\), the new matrix is
\[
        (U_-L)^*(U_+R)=L^*U_-^*U_+R=L^*AR.
\]
All orthonormal frames of a fixed finite-dimensional subspace differ by a
unitary on the coordinate space, so the converse follows. The rectangular SVD
classifies the biunitary orbit by singular values, including rank-deficient
zeros and rectangular padding.
\end{proof}

\begin{corollary}[No canonical entrywise matrix without frames]
From a fixed subspace pair alone, a sorted diagonal SVD representative is an
orbit normal form, not an endpoint-selected matrix. Any entrywise matrix
requires an additional frame-selection rule.
\end{corollary}

\subsection{When principal frames exist}

\begin{theorem}[Simple-spectrum frame criterion]
On the nonzero-overlap sector, canonical matched principal frames can be
selected from the subspace pair alone exactly on simple nonzero principal
spectra. If a nonzero principal value has multiplicity \(m\ge2\), no natural
orthonormal principal frame is determined by the subspace pair alone.
\end{theorem}

\begin{proof}
If the nonzero eigenvalues of \(K_+\) are simple, their eigenspaces are complex
lines. Choosing unit vectors in these lines is unique up to phase. Applying the
polar partial isometry gives the matched principal lines in \(S_-\). Thus
principal frames are canonical up to independent phases.

If a nonzero eigenvalue has multiplicity \(m\ge2\), the corresponding
eigenspace carries an internal \(U(m)\) freedom. In the model block where
\(C=\sigma\Id\) on an \(m\)-dimensional principal block, every orthonormal basis
is equally compatible with the same subspace-pair data. A natural frame rule
depending only on the subspace pair would have to be fixed by all of \(U(m)\),
which is impossible for an ordered orthonormal basis when \(m\ge2\).
\end{proof}

\begin{corollary}[Finite-dimensional completeness]
For finite-dimensional subspace pairs, the endpoint dimensions and the
principal-angle multiset classify the pair up to ambient unitary equivalence.
\end{corollary}

\begin{proof}
This is the finite-dimensional CS decomposition. The pair splits into common,
orthogonal, residual, and principal two-plane blocks. The dimensions and
principal angles determine the block model, and any pair with the same data is
unitarily equivalent to that model.
\end{proof}

\begin{gap}[GAP-5A: theorem packaging]
The endpoint operator package is now mathematically stable, but the final paper
should package it as one clean theorem with named corollaries: projector test,
polar transport, frame criterion, biunitary finality, and CS completeness. This
is mainly an exposition gap, not a proof gap.
\end{gap}

\section{\texorpdfstring{\(A_r\) versus \(O_r\)}{A r versus O r}}

Two related subspace filtrations occur throughout the paper:
\[
        O_r=\Span\{P\partial^\alpha\psi:1\le|\alpha|\le r\},
\]
and
\[
        A_r=\Span\{\psi,\partial^\alpha\psi:1\le|\alpha|\le r\}.
\]
The first is horizontal and value-channel-free. The second keeps the state line
and is often the cleaner ambient object.

\begin{proposition}[Role split]
For each endpoint,
\[
        A_r=\Span\{\psi\}\oplus O_r.
\]
The orthogonal projector \(P=\Id-|\psi\rangle\langle\psi|\) restricts to a
surjection \(A_r\to O_r\) with kernel \(\Span\{\psi\}\). Thus \(O_r\) is
canonically the quotient of \(A_r\) by the state line.
\end{proposition}

\begin{proof}
Each derivative decomposes into its horizontal projection plus a state-line
component. This gives \(A_r\subseteq \Span\{\psi\}+O_r\), while the reverse
inclusion is immediate from the definition of \(O_r\). Orthogonality of
\(\Span\{\psi\}\) and \(O_r\) follows from the horizontal projection. The
kernel and quotient statement follow from the restriction of \(P\).
\end{proof}

There is no global winner between these two objects. \(O_r\) is the faithful
horizontal object, closest to the value-channel-removed QGT pattern. \(A_r\) is
the ambient extension and supports the cleanest transport-free endpoint theorem
when all endpoints live in a common Hilbert space. Applications should state
which role is being used.

\begin{gap}[GAP-6A: physical selection principle]
The current theorem is a role split, not a selection principle. A stronger
physics-facing paper would identify settings where \(A_r\) or \(O_r\) is forced
by an operational or dynamical criterion. For now, the paper keeps both and
states the distinction explicitly.
\end{gap}

\section{Dynamics and operational readout}

\subsection{Exact unitary orbits and Krylov spaces}

Let \(\psi(t)=e^{-itH}\psi_0\). Then
\(\partial_t^k\psi(t)=(-iH)^ke^{-itH}\psi_0\), so
\[
        A_r(t)=e^{-itH}\Span\{\psi_0,H\psi_0,\ldots,H^r\psi_0\}.
\]
The horizontal space \(O_r(t)\) is the value-channel-removed counterpart:
\[
        O_r(t)=e^{-itH}\bigl(\Span\{\psi_0,H\psi_0,\ldots,H^r\psi_0\}
                 \cap \psi_0^\perp\bigr).
\]
Thus exact unitary dynamics turns the abstract osculating spaces into
transported Krylov and compressed-Krylov spaces.

If the cyclic space of \((H,\psi_0)\) has dimension \(m\), then the spaces
stabilize once \(r\ge m-1\). In a finite-cyclic case, the spectral measure of
\((H,\psi_0)\) gives an orthogonal-polynomial basis and a finite Jacobi matrix
\(J\). In that basis, the full ambient unitary package is represented by
\(e^{-i\Delta J}\). This is a specialization of the subspace theorem, not a
general matrix theorem.

\subsection{Projector speed and leakage}

For a unitary subspace orbit \(S(t)=e^{-itH}S(0)\), let \(\Proj_t\) be the
orthogonal projector onto \(S(t)\). Then
\[
        \dot\Proj_t=-i[H,\Proj_t]
\]
and
\[
        \frac12\|\dot\Proj_t\|_{HS}^2
        =
        \frac12\|[H,\Proj_t]\|_{HS}^2
        =
        \|(\Id-\Proj_t)H\Proj_t\|_{HS}^2 .
\]
The Grassmannian speed is the Hilbert-Schmidt size of the Hamiltonian leakage
block. At finite separation,
\[
        \Proj_t\Proj_{t+\Delta}\Proj_t|_{S(t)}
        =
        \Id-\Delta^2L_t^*L_t+O(\Delta^3),
        \qquad L_t=(\Id-\Proj_t)H\Proj_t.
\]

\subsection{Kato transport and curvature}

Every smooth constant-rank projector path admits the Kato Hamiltonian
\[
        H_K(t)=i[\dot\Proj(t),\Proj(t)].
\]
Its unitary propagator realizes the path:
\[
        \Proj(t)=U_K(t)\Proj(0)U_K(t)^*.
\]
This gives a canonical pathwise model. It is not an endpoint-only comparison.
In several parameters the Kato connection has curvature
\[
        R^K=\Proj(\dd\Proj\wedge\dd\Proj)\Proj.
\]
Nonzero curvature produces holonomy and obstructs endpoint-independent Kato
transport. In rank one this reduces to Berry curvature.

\subsection{Reflection protocol}

The endpoint operator package also has a coherent readout under an explicit
access model.

\begin{proposition}[Alternating reflections read principal angles]
Let \(S_\pm\) be finite-dimensional endpoint subspaces with projectors
\(\Proj_\pm\). Suppose the reflections
\[
        R_\pm=2\Proj_\pm-\Id
\]
can be implemented coherently. On each nontrivial principal two-plane of the
pair, the unitary \(W=R_-R_+\) has eigenphases \(\pm2\theta_i\), where
\(\theta_i\) is the corresponding principal angle.
\end{proposition}

\begin{proof}
On a principal two-plane choose an orthonormal basis \(e,f\) with
\[
        e\in S_-,\qquad
        S_+\cap\Span\{e,f\}
        =
        \Span\{\cos\theta\,e+\sin\theta\,f\}.
\]
In this basis
\[
        R_-=
        \begin{pmatrix}
        1&0\\0&-1
        \end{pmatrix},
        \qquad
        R_+=
        \begin{pmatrix}
        \cos2\theta&\sin2\theta\\
        \sin2\theta&-\cos2\theta
        \end{pmatrix}.
\]
Thus
\[
        R_-R_+=
        \begin{pmatrix}
        \cos2\theta&\sin2\theta\\
        -\sin2\theta&\cos2\theta
        \end{pmatrix},
\]
whose eigenvalues are \(e^{\pm 2i\theta}\), up to the convention for rotation
orientation.
\end{proof}

Intersection sectors give phase \(+1\), one-sided orthogonal residual sectors
give phase \(-1\), and ambient padding must be separated from the endpoint
principal-angle data. Controlled phase estimation on \(W\) estimates the
principal angles when coherent reflection access and adequate state preparation
or sampling are available. This is a protocol for reading the existing
operator package. It is not a new invariant and not a universal Hamiltonian
model for arbitrary jet families.

\begin{gap}[GAP-7A: beyond reflection access]
The reflection protocol is natural in quantum algorithmic settings with
coherent membership tests, code-space reflections, symmetry-sector reflections,
or phase-kickback access. It may be too strong for arbitrary jet-derived
subspaces. A stronger physics paper should find a concrete Hamiltonian,
measurement, or response model where the same principal angles arise without
assuming arbitrary projector reflections.
\end{gap}

\section{Benchmark families}

The theory is not only a no-go theorem. The surviving subspace data can be
richer than endpoint overlap.

\subsection{Oscillator sanity check}

For the harmonic-oscillator ground state
\[
        \psi_\omega(x)=
        \left(\frac{\omega}{\pi}\right)^{1/4}e^{-\omega x^2/2},
        \qquad \omega>0,
\]
the same-point first-jet block in the \(\omega\)-coordinate is
\[
        M(\omega)=
        \begin{pmatrix}
        1&0\\
        0&\frac{1}{8\omega^2}
        \end{pmatrix}.
\]
The endpoint overlap is
\[
        I(\omega_-,\omega_+)=
        \frac{\sqrt2(\omega_-\omega_+)^{1/4}}
             {\sqrt{\omega_-+\omega_+}}.
\]
The raw first-jet blocks change under the squeeze parameter
\(r=\frac12\log\omega\), but the bilaterally whitened first-jet block stays
unchanged in this benchmark. Its singular values collapse to a function of
overlap alone, so the oscillator is a sanity check rather than a
richer-than-overlap example.

\subsection{Qutrit polynomial curve}

Consider
\[
        \psi(t)=\frac{(1,t,t^2)^T}{\sqrt{1+t^2+t^4}}.
\]
Writing
\[
        D(t)=1+t^2+t^4,\qquad J(t)=1+4t^2+t^4,
\]
the identity-transport representative of the first-horizontal whitened block is
\[
\Omega(u,v)=
\begin{pmatrix}
\frac{1+uv+u^2v^2}{\sqrt{D(u)D(v)}} &
\frac{(u-v)(1+2v^2+2uv+uv^3)}{\sqrt{D(u)D(v)J(v)}}\\[1.2ex]
\frac{(v-u)(1+2u^2+2uv+u^3v)}{\sqrt{D(u)D(v)J(u)}} &
\frac{1-u^4-v^4+5uv+4u^3v+4uv^3+5u^3v^3+u^4v^4}
{\sqrt{D(u)D(v)J(u)J(v)}}
\end{pmatrix}.
\]
Its singular values are \(1\) and
\[
        \frac{|1+4uv+u^2v^2|}
        {\sqrt{(1+4u^2+u^4)(1+4v^2+v^4)}}.
\]
The pairs
\[
        (u,v)=\left(-5,\frac65\right),
        \qquad
        (u',v')=\left(-\frac52,-\frac35\right)
\]
have the same overlap
\[
        I(u,v)=I(u',v')=\frac{25\sqrt{39}}{273},
\]
but different second singular values. Thus the first-order endpoint geometry
sees more than ordinary overlap in this benchmark.

\subsection{\texorpdfstring{Quartit cubic and \(O_2\)}{Quartit cubic and O2}}

For
\[
        \psi(t)=\frac{(1,t,t^2,t^3)^T}
        {\sqrt{1+t^2+t^4+t^6}},
\]
the ambient second osculating space is a varying \(3\)-plane in \(\R^4\), with
normal
\[
        n(t)=(-t^3,3t^2,-3t,1)^T.
\]
The surviving ambient singular values are
\[
        1,\qquad 1,\qquad
        \frac{|1+9uv+9u^2v^2+u^3v^3|}
        {\sqrt{(1+9u^2+9u^4+u^6)(1+9v^2+9v^4+v^6)}}.
\]
There are endpoint pairs with the same overlap and different final singular
value.

The value-channel-free side is more informative. For
\[
        O_2(t)=\Span\{j_1(t),j_2(t)\},
\]
the principal-angle cosines of \(O_2(u)\) and \(O_2(v)\) are the singular
values of
\[
M_{O_2}(u,v)=
\begin{pmatrix}
\frac{1+uv+u^2v^2+u^3v^3}{\sqrt{D_4(u)D_4(v)}} &
\frac{(u-v)^3}{\sqrt{D_4(u)N_4(v)}}\\[1.2ex]
-\frac{(u-v)^3}{\sqrt{N_4(u)D_4(v)}} &
\frac{1+9uv+9u^2v^2+u^3v^3}{\sqrt{N_4(u)N_4(v)}}
\end{pmatrix},
\]
where
\[
        D_4(t)=1+t^2+t^4+t^6,\qquad
        N_4(t)=1+9t^2+9t^4+t^6.
\]
On this benchmark \(O_2\) is strictly finer than \(A_2\): the ambient invariant
keeps only the product of the two \(O_2\) singular values, while \(O_2\) keeps
both.

\subsection{Veronese family}

Let
\[
        v_n(t)=(1,t,\ldots,t^n).
\]
At top order \(r=n-1\), the ambient osculating subspace is a codimension-one
hyperplane. Its normal is the coefficient vector of \((x-t)^n\). Therefore the
single nontrivial principal-angle cosine is the normalized binomial-square
polynomial in \(uv\). For every \(n\ge2\), this quantity is generically not a
function of endpoint overlap alone.

On the value-channel-free side, \(O_{n-1}\) has a codimension-two complement
description. The corresponding \(2\times2\) complement matrix recovers the
qutrit \(O_1\) and quartit \(O_2\) cases and is generically richer than overlap.
Lower-order \(r<n-1\) complements are known structurally, but closed spectral
formulas are not yet part of the paper.

\subsection{Spherical twist family}

The phenomenon is not confined to polynomial curves. For
\[
        \psi_k(t)=(\cos t\cos kt,\cos t\sin kt,\sin t),
        \qquad k\ge2,
\]
there are exact equal-overlap endpoint pairs with different \(A_1\) and \(O_1\)
principal-angle data. For \(k=2\), the pairs \((0,\pi/4)\) and \((0,\pi/2)\)
both have overlap \(0\), but the corresponding \(A_1\) and \(O_1\) angle data
differ.

\subsection{\texorpdfstring{Finite-jet genericity for \(O_1\)}{Finite-jet genericity for O1}}

The examples above have a local generic mechanism. For normalized real \(C^2\)
curves in finite dimension at least \(3\), fix two distinct parameter values
with nonzero endpoint velocities. Define
\[
        I(u,v)=\langle\psi(u),\psi(v)\rangle,\qquad
        T(t)=\frac{\psi'(t)}{\|\psi'(t)\|},\qquad
        C(u,v)=\langle T(u),T(v)\rangle .
\]
On a branch where \(C\ne0\), the first value-channel-free line-angle observable
is locally independent of overlap whenever
\[
        \det\frac{\partial(I,C)}{\partial(u,v)}\ne0.
\]
After clearing the nonzero velocity denominators, this determinant is a
polynomial in the endpoint \(2\)-jets and is not identically zero on the
admissible jet chart. Thus nonvanishing is an open dense finite-jet condition.
Consequently, for smooth curve families whose endpoint \(2\)-jet map is
submersive onto such a chart, and for real-analytic families where the pulled
back determinant is not identically zero, there is an open dense locus where the
\(O_1\) line-angle data is not locally determined by overlap alone.

\begin{gap}[GAP-8A: higher-order genericity]
The finite-jet theorem is first-order and value-channel-free. A stronger paper
would prove generic richer-than-overlap behavior for higher \(O_r\) and \(A_r\)
principal-angle spectra, and would identify the sharp dimension and regularity
hypotheses.
\end{gap}

\begin{gap}[GAP-8B: lower-order Veronese spectra]
The lower-order Veronese complement spaces have structural descriptions, but
closed or conceptual formulas for their full principal-angle spectra remain to
be written. This is a concrete algebraic target for the next research cycle.
\end{gap}

\section{Limits and research program}

The theory has a clean core and clear boundaries.

\begin{enumerate}
\item It does not produce a canonical full matrix from higher jets. It explains
why such a matrix is absent without extra frame-selecting structure.
\item It does not give endpoint-independent Kato transport in nonzero
curvature. Kato transport is pathwise.
\item It does not choose principal frames in repeated-spectrum sectors.
\item It does not yet supply a physics case study where higher osculating
subspace angles reveal a new material, algorithmic, or dynamical effect beyond
known QGT/fidelity diagnostics.
\end{enumerate}

The next research program is therefore concrete.

\begin{gap}[GAP-10A: physics case study]
Find a concrete quantum system where the \(O_r/A_r\) principal-angle package
reveals information not visible to ordinary overlap, fidelity susceptibility,
or the first-order quantum geometric tensor. Candidate settings include Bloch
state families, flat-band geometry, topological phase boundaries, variational
ansatz manifolds, and Krylov dynamics.
\end{gap}

\begin{gap}[GAP-10B: mixed states and degenerate bands]
Extend, or sharply fail to extend, the osculating-subspace package to density
matrices, purification bundles, or degenerate Bloch bands. This must be compared
with non-Abelian quantum geometry and Wilczek-Zee transport before any novelty
claim is made.
\end{gap}

\begin{gap}[GAP-10C: numerical invariant library]
Build a stable numerical pipeline that takes a sampled state family, estimates
the \(O_r\) and \(A_r\) endpoint subspaces, and computes principal-angle spectra
with conditioning diagnostics. This would turn the theory into a usable
diagnostic tool.
\end{gap}

\section{Conclusion}

The attempted higher-order matrix generalization of the quantum geometric
tensor fails for principled reasons. Higher projected jets mix by non-unitary
triangular laws under phase and coordinate changes, so raw matrix entries are
not canonical. The failure identifies the correct invariant object: the
osculating subspace filtration and the endpoint principal-angle/operator
package.

This object is mathematically rigid. Its frame matrices form a biunitary orbit,
its singular values are principal-angle cosines, its polar factor is canonical
on the nonzero-overlap sector, and its finite-dimensional classification is the
CS decomposition. It is also useful: explicit benchmarks show it can distinguish
endpoint pairs with the same ordinary overlap, and coherent alternating
reflections read it under standard access assumptions.

The remaining work is not to force a noncanonical matrix into the theory. It is
to use the canonical subspace/operator package in serious physics examples.

\section*{Provenance note}

This draft is assembled from the working note, the paper skeleton, and the
reports in the active quantum team directory. The most important deposits are
the reports on the first-order QGT bridge, the higher-jet no-go, covariant jets,
the endpoint operator package, biunitary matrix finality, the reflection
protocol, and nonbenchmark genericity.

\end{document}
